\chapter{Modélisation des produits à taux overnight}\label{ch:rc:modelling-overnight-rate-products}

Un prêt à taux variable d'une entreprise verse, chaque trimestre, le SOFR composé quotidiennement sur
le trimestre, connu seulement à la fin de celui-ci. Sa trésorière achète un cap, et
le desk qui le vend remarque quelque chose que les anciens caps interbancaires n'avaient jamais fait~: le
premier caplet, dont la période a déjà commencé, conserve encore une optionnalité, et
chaque caplet continue d'accumuler du risque après le début de sa période, jusqu'à son dernier
fixing overnight. Un taux interbancaire est fixé au début de sa période~; un
taux overnight composé se fixe à la fin. Les modèles des chapitres~7 et~8 ont été
construits pour des taux connus à l'avance. Ce chapitre les étend à des taux connus
a posteriori~: dans le \omterm{def:rc:short-rate-models:hw}{modèle de Hull--White}, le \omterm{def:rc:modelling-overnight-rate-products:caplet}{caplet rétrospectif} admet une
forme fermée~; dans le modèle de marché, l'extension est le \omterm{def:rc:modelling-overnight-rate-products:gfmm}{modèle de marché forward généralisé} (FMM), et
avec le calendrier des réunions des banques centrales, le taux devient une fonction en escalier
dont le risque se concentre aux jours de réunion.

\section{Taux rétrospectifs et leurs options}

\begin{definition}[Taux rétrospectif, taux à terme prospectif]\label{def:rc:modelling-overnight-rate-products:rates}
Un \emph{taux rétrospectif}\index{taux rétrospectif} pour une période $[S,E]$ est
le taux overnight composé a posteriori sur la période (livre~2,
chapitre~1)~: $1+\delta R(S,E) = \prod_i(1+r_id_i/360)$, connu seulement en $E$. Un
\emph{taux à terme prospectif}\index{taux à terme prospectif} pour la même
période est publié en $S$, comme espérance de ce taux composé implicite dans
les contrats à terme et les swaps (un taux SOFR à terme) ou comme taux interbancaire offert (Euribor)~;
il est connu dès le début.
\end{definition}

\begin{definition}[Caplet rétrospectif]\label{def:rc:modelling-overnight-rate-products:caplet}
Un \emph{caplet rétrospectif}\index{caplet rétrospectif} verse
$\delta\max(R(S,E)-K,0)$ en $E$ (ou quelques jours plus tard) sur le taux composé de
sa propre période.
\end{definition}

\begin{omfigure}
\begin{tikzpicture}[x=2.0cm,font=\small]
\draw[->,thick] (0,0) -- (5.4,0) node[right] {temps};
\foreach \x/\l in {0/0,2/$S$,4/$E$} {\draw (\x,-0.1) -- (\x,0.1); \node[below] at (\x,-0.12) {\l};}
\draw[omDef,thick,fill=omDef!12] (2,0.4) rectangle (4,0.8); \node[font=\footnotesize] at (3,0.6) {période de calcul};
\draw[omThm,very thick] (0,1.4) -- (2,1.4); \draw[omThm,thick,fill=omThm] (2,1.4) circle (2pt);
\node[text=omThm,anchor=west,font=\footnotesize] at (4.15,1.4) {prospectif~: risque jusqu'à $S$, taux connu en $S$};
\draw[omProp,very thick] (0,2.1) -- (4,2.1); \draw[omProp,thick,fill=omProp] (4,2.1) circle (2pt);
\node[text=omProp,anchor=west,font=\footnotesize] at (4.15,2.1) {rétrospectif~: risque jusqu'à $E$, taux connu en $E$};
\end{tikzpicture}

\omcaption{Un taux prospectif porte un risque jusqu'au début de sa période
et est connu dès lors~; un taux overnight composé continue de varier jusqu'au
dernier fixing, en fin de période, si bien qu'une option sur lui porte un risque sur toute
sa période de calcul.}\label{fig:rc:modelling-overnight-rate-products:timing}
\end{omfigure}

\begin{proposition}[Caplets rétrospectifs dans Hull--White]\label{prop:rc:modelling-overnight-rate-products:hw}
Approximons la capitalisation quotidienne par $1+\delta R = \exp\int_S^Er_u\,du$. Dans le \omterm{def:rc:short-rate-models:hw}{modèle de Hull--White}
avec $\sigma$ constant et $c = 1+\delta K$,
\[
\mathrm{Caplet}_0 = P(0,S)\bigl[\Phi(-d_2) - c\,m\,\Phi(-d_1)\bigr],\quad
m = \frac{P(0,E)}{P(0,S)},\ d_1 = \frac{\ln(mc)+V/2}{\sqrt V},\ d_2 = d_1-\sqrt V,
\]
avec $V = V_{\text{avant}}+V_{\text{pendant}}$,
$V_{\text{avant}} = \sigma^2B(\delta)^2\frac{1-e^{-2\kappa S}}{2\kappa}$ et
$V_{\text{pendant}} = \sigma^2\int_0^\delta B(\tau)^2\,d\tau\approx\sigma^2\delta^3/3$. Le
caplet à terme prospectif sur un taux à terme pour la même période a la même formule
avec $V = V_{\text{avant}}$ seulement.
\end{proposition}
\begin{proof}
Le gain vaut $(e^{-\int_0^Sr}-c\,e^{-\int_0^Er})^+$. Sous la mesure de densité
$e^{-\int_0^Sr}/P(0,S)$, le prix vaut $P(0,S)\E[(1-cL)^+]$ avec
$L=e^{-\int_S^Er}$, lognormale de moyenne $P(0,E)/P(0,S)$~; sa variance logarithmique est la
variance de $\int_S^Ex_u\,du$, qui se décompose en chocs avant $S$ (à travers le
prix du zéro-coupon en $S$) et ceux pendant la période. L'espérance est un put de Black.
\end{proof}

\omcode{code/firm/rfrcaplet/firm_rfrcaplet.py}{35}{56}{Les deux parties de la variance, et le caplet rétrospectif sous forme fermée.}\label{lst:rc:modelling-overnight-rate-products:hw}

\begin{example}[Un cap SOFR de deux ans]\label{ex:rc:modelling-overnight-rate-products:cap}
Sur la courbe SOFR du chapitre~1 avec $\kappa=3~\%$ et $\sigma=90$ points de base, un cap de deux ans sur
le SOFR composé à trois mois à 3.75~\% sur 100~millions de dollars US coûte 288\,963~dollars US~: huit
caplets, de 13\,623 (le trimestre en cours) à 53\,846
(voir la \cref{fig:rc:modelling-overnight-rate-products:caplets}). Le même cap sur un
taux à terme à trois mois, dont le premier caplet a déjà fixé, coûte 250\,741~dollars US~:
le cap rétrospectif coûte 38\,222~dollars US de plus, dont 13\,623 pour le
premier caplet et 24\,599 pour la variance pendant la période des sept autres. Une simulation de Monte Carlo
du taux court avec 260 pas par an réévalue le caplet du cinquième trimestre
(40\,944~dollars US) à moins d'un cinquième d'erreur type près.
\end{example}

\begin{omfigure}
\begin{tikzpicture}
\begin{axis}[width=12cm,height=5.4cm,ybar,bar width=10pt,
  ylabel={valeur du caplet (milliers de dollars US)},xlabel={période (années depuis le spot)},
  tick label style={font=\small},label style={font=\small},
  xmin=-0.6,xmax=7.6,ymin=0,ymax=60,grid=major,grid style={gray!25},
  xtick=data,xticklabels from table={figdata/rates-credit-risk/10-modelling-overnight-rate-products/caplets.csv}{label},
  x tick label style={font=\scriptsize},
  legend columns=2,legend style={font=\footnotesize,at={(0.5,-0.3)},anchor=north,draw=none,fill=none,/tikz/every even column/.append style={column sep=8pt}},
  table/col sep=comma]
\addplot[fill=omProp!55,draw=omProp] table[x=i,y=back]{figdata/rates-credit-risk/10-modelling-overnight-rate-products/caplets.csv};
\addplot[fill=omDef!55,draw=omDef] table[x=i,y=fwd]{figdata/rates-credit-risk/10-modelling-overnight-rate-products/caplets.csv};
\legend{rétrospectif (SOFR composé),prospectif (taux à terme)}
\end{axis}
\end{tikzpicture}

\omcaption{Caplets d'un cap de deux ans à 3.75~\% sur 100~millions de dollars US, sur le SOFR composé
et sur un taux à terme pour les mêmes trimestres. Le trimestre en cours n'a de caplet
que dans le cap rétrospectif~; chaque \omterm{def:rc:modelling-overnight-rate-products:caplet}{caplet rétrospectif} ultérieur vaut davantage
par la variance a posteriori de sa période. Données~: courbe du chapitre~1, Hull--White $\kappa=3~\%$,
$\sigma=90$ points de base~; le tutoriel du chapitre.}\label{fig:rc:modelling-overnight-rate-products:caplets}
\end{omfigure}

La part de variance pendant la période vaut 100~\% pour le trimestre en cours, 25~\% pour
le suivant, 7.9~\% pour le cinquième et 4.8~\% pour le huitième~: proche de $(\delta/3)/(S+\delta/3)$,
le rapport que donne une décroissance linéaire de la volatilité sur la période
(voir la \cref{fig:rc:modelling-overnight-rate-products:share}).

\begin{omfigure}
\begin{tikzpicture}
\begin{axis}[width=11cm,height=5cm,
  xlabel={début du trimestre $S$ (années)},ylabel={part de la variance pendant la période (\%)},
  tick label style={font=\small},label style={font=\small},
  xmin=0,xmax=1.8,ymin=0,ymax=105,grid=major,grid style={gray!25},
  legend columns=2,legend style={font=\footnotesize,at={(0.5,-0.25)},anchor=north,draw=none,fill=none,/tikz/every even column/.append style={column sep=8pt}},
  table/col sep=comma]
\addplot[omThm,thick,mark=*] table[x=s,y=hw]{figdata/rates-credit-risk/10-modelling-overnight-rate-products/share.csv};
\addplot[omDef,thick,dashed] table[x=s,y=lin]{figdata/rates-credit-risk/10-modelling-overnight-rate-products/share.csv};
\legend{Hull--White,décroissance linéaire~: $(\delta/3)/(S+\delta/3)$}
\end{axis}
\end{tikzpicture}

\omcaption{Part de la variance d'un \omterm{def:rc:modelling-overnight-rate-products:caplet}{caplet rétrospectif} qui s'accumule pendant son propre
selon le début du trimestre~: la totalité pour le trimestre en cours, un quart pour le
suivant, quelques pour cent à deux ans. Données~: le
tutoriel du chapitre.}\label{fig:rc:modelling-overnight-rate-products:share}
\end{omfigure}

\section{Le modèle de marché forward généralisé}

\begin{definition}[Modèle de marché forward généralisé]\label{def:rc:modelling-overnight-rate-products:gfmm}
Le \emph{modèle de marché forward généralisé}\index{modèle de marché forward généralisé}
(Lyashenko et Mercurio, 2019) modélise, pour chaque période de calcul $[T_{k-1},T_k]$,
le forward $R_k(t) = \E^{T_k}_t[R(T_{k-1},T_k)]$ du \omterm{def:rc:modelling-overnight-rate-products:rates}{taux rétrospectif},
pour tout $t\le T_k$~: avant $T_{k-1}$, il coïncide avec le forward prospectif
du chapitre~8, et pendant la période il continue de suivre les fixings au jour le jour. Sa
volatilité vaut $\sigma_k(t)g_k(t)$, avec $g_k=1$ avant $T_{k-1}$ et $g_k$ décroissant
vers zéro en $T_k$, par exemple linéairement.
\end{definition}

\begin{proposition}[Variance effective avec décroissance linéaire]\label{prop:rc:modelling-overnight-rate-products:gfmm}
Avec $\sigma_k$ constant et $g_k(t) = (T_k-t)/(T_k-T_{k-1})$ sur la période, le
\omterm{def:rc:modelling-overnight-rate-products:caplet}{caplet rétrospectif} se valorise par la formule de Black avec une variance totale
$\sigma_k^2\bigl(T_{k-1}+\tfrac13(T_k-T_{k-1})\bigr)$, contre $\sigma_k^2T_{k-1}$ pour le
caplet prospectif.
\end{proposition}
\begin{proof}
$\int_0^{T_{k-1}}\sigma_k^2\,dt+\int_{T_{k-1}}^{T_k}\sigma_k^2\bigl(\frac{T_k-t}{\delta}\bigr)^2dt =
\sigma_k^2(T_{k-1}+\delta/3)$, et $R_k$ est lognormal sous $\mathbb Q^{T_k}$.
\end{proof}

Le modèle généralisé contient l'ancien~: un même processus donne à la fois le
taux prospectif et le \omterm{def:rc:modelling-overnight-rate-products:rates}{taux rétrospectif}, si bien qu'un portefeuille de produits interbancaires historiques
et de produits au jour le jour nouveaux est géré du point de vue du risque dans un seul modèle, avec la
l'outillage de simulation et de calibration du modèle de marché.

\section{Modèles markoviens de faible dimension}

Les modèles de marché sont de grande dimension~; les systèmes de production pour les exotiques et
les expositions privilégient des modèles à une poignée de variables d'état dont la courbe à terme
admet une formule exacte.

\begin{definition}[Modèle de Cheyette]\label{def:rc:modelling-overnight-rate-products:cheyette}
Un \emph{modèle de Cheyette}\index{modèle de Cheyette} (1992) est un modèle HJM dont la volatilité forward
se factorise en $\sigma_f(t,T) = \sigma_r(t,x_t,y_t)\,e^{-\kappa(T-t)}$. Alors
$f(t,T) = f(0,T)+e^{-\kappa(T-t)}\bigl(x_t+B(t,T)y_t\bigr)$ avec deux variables d'état,
$dx = (y-\kappa x)\,dt+\sigma_r\,dW$ et $dy = (\sigma_r^2-2\kappa y)\,dt$. Avec
$\sigma_r$ déterministe, c'est Hull--White~; laisser $\sigma_r$ dépendre de $x$ donne un
\omterm{def:rc:short-rate-models:short}{modèle de taux court} à volatilité locale (avec skew) qui reste markovien en $(x,y)$.
\end{definition}

L'intégrale du taux court sur une période de calcul est une variable d'état de plus
dans un tel modèle, si bien que les produits rétrospectifs restent de faible dimension~: c'est la
voie markovienne vers les mêmes caplets, celle que suivent les méthodes de pricing par EDP et par arbres.

\section{Dates de réunion et sauts du taux court}

Le taux au jour le jour ne diffuse pas~: il reste au taux directeur et ne bouge que les
jours où une banque centrale décide (le livre~2, chapitre~8). Son incertitude
arrive par à-coups, selon un calendrier publié.

\begin{dated}{2026-09}{Calendrier des réunions et taux à terme}\label{dat:rc:modelling-overnight-rate-products:term}
Le comité de politique monétaire de la Federal Reserve se réunit les 27--28 octobre et 8--9 décembre
2026, puis les 26--27 janvier, 16--17 mars, 27--28 avril, 8--9 juin, 27--28 juillet,
14--15 septembre, 26--27 octobre et 7--8 décembre 2027 (provisoire jusqu'à confirmation
lors de la réunion précédente). L'ARRC a formellement recommandé les taux SOFR à terme prospectifs du CME Group
le 29 juillet 2021, et ses bonnes pratiques ne soutiennent pas leur usage dans les produits
dérivés, sauf chez les utilisateurs finaux qui couvrent des produits au comptant qui y font référence.
\end{dated}

\omcode{code/firm/rfrcaplet/firm_rfrcaplet.py}{106}{111}{La variance de la moyenne du taux sur une période lorsque le taux ne bouge qu'aux réunions.}\label{lst:rc:modelling-overnight-rate-products:meetings}

\begin{example}[Un trimestre, deux réunions]\label{ex:rc:modelling-overnight-rate-products:meetings}
Le trimestre du 29 septembre au 29 décembre 2027 (années 1.0 à 1.25 à partir du spot)
contient les réunions des 27 octobre et 8 décembre. Si le taux ne bouge qu'aux
réunions, une décision du 27 octobre déplace la moyenne du taux du trimestre de 69~\% du
saut (la part du trimestre restante), une du 8 décembre de 23~\%~; après le 8 décembre le
taux du trimestre est connu, trois semaines avant la fin. Le modèle généralisé à
décroissance linéaire répartit le même risque en douceur
(\cref{fig:rc:modelling-overnight-rate-products:profile})~: ramenés à un même total, les deux
coïncident avant le trimestre et diffèrent à l'intérieur de celui-ci.
\end{example}

\begin{omfigure}
\begin{tikzpicture}
\begin{axis}[width=12cm,height=5.4cm,
  xlabel={temps $t$ (années depuis le spot)},ylabel={écart type restant (\% du total)},
  tick label style={font=\small},label style={font=\small},
  xmin=0,xmax=1.25,ymin=0,ymax=105,grid=major,grid style={gray!25},
  legend columns=2,legend style={font=\footnotesize,at={(0.5,-0.25)},anchor=north,draw=none,fill=none,/tikz/every even column/.append style={column sep=8pt}},
  table/col sep=comma]
\addplot[omDef,very thick] table[x=t,y=gfmm]{figdata/rates-credit-risk/10-modelling-overnight-rate-products/profile.csv};
\addplot[omThm,thick,const plot] table[x=t,y=meet]{figdata/rates-credit-risk/10-modelling-overnight-rate-products/profile.csv};
\legend{modèle de marché forward généralisé (décroissance linéaire),taux ne bougeant qu'aux réunions}
\draw[gray,dashed] (axis cs:1.0,0) -- (axis cs:1.0,105) node[pos=0.88,left,font=\footnotesize,text=black] {début de la période};
\end{axis}
\end{tikzpicture}

\omcaption{Écart type restant, vu en $t$, du taux composé du
trimestre de l'année 1.0 à 1.25, dans deux modèles ramenés au même total~: une décroissance
régulière sur toute la période, ou des paliers aux jours de réunion de la Federal Reserve (dix
réunions avant la fin du trimestre~; la dernière à l'intérieur de celui-ci le 8 décembre 2027). Données~:
calendrier des réunions de la Federal Reserve~; le
tutoriel du chapitre.}\label{fig:rc:modelling-overnight-rate-products:profile}
\end{omfigure}

\begin{remark}[Ce que change la vue par réunion]\label{rem:rc:modelling-overnight-rate-products:meetings}
Les prix changent peu, puisque la variance totale est calibrée dans les deux cas~; le
risque change beaucoup. Le theta et le vega d'un caplet se concentrent aux jours de réunion, et un
portefeuille d'options overnight à court terme doit être couvert réunion par réunion, avec
les contrats à terme du livre~2, chapitre~8, et les options sur ceux-ci.
\end{remark}

\section{Taux à terme prospectifs}

Les taux à terme subsistent pour les prêts~: un emprunteur veut connaître les
intérêts du trimestre suivant dès le départ. Dérivé de contrats à terme et de swaps overnight, un taux
SOFR à terme est une anticipation publiée~; les options et les swaps portant dessus sont des
instruments prospectifs de la nature de ceux du chapitre~4. Le marché garde les dérivés sur le taux overnight
lui-même, là où se trouve la liquidité, et réserve les dérivés de taux à terme aux couvertures des
prêts qui les utilisent, afin que le taux à terme ne repose pas sur le trading
portant sur lui-même. Le desk qui vend un cap sur taux à terme à un emprunteur le couvre donc
avec des options sur taux overnight et porte la base entre les deux~: la
variance intra-période de l'\cref{ex:rc:modelling-overnight-rate-products:cap}.

\section{Tutoriel~: rétrospectif contre prospectif}\label{tut:rc:modelling-overnight-rate-products:tutorial}

\begin{tutorial}
\textbf{Objectif.} Évaluer un cap à deux ans sur le SOFR composé et sur un taux à terme, vérifier la
formule fermée par simulation, et comparer les profils de risque régulier et ceux aux dates de réunion.
\textbf{État final~:}
les \cref{fig:rc:modelling-overnight-rate-products:caplets,fig:rc:modelling-overnight-rate-products:profile}
et les chiffres de l'\cref{ex:rc:modelling-overnight-rate-products:cap}.
\begin{tutsteps}
\item \textbf{Variances}~: \texttt{hw\_period\_variances(kappa, sigma, S, E)} pour chaque
trimestre.
\item \textbf{Caps}~: \texttt{cap\_table()} et \texttt{cap\_totals()}.
\item \textbf{Simuler}~: \texttt{mc\_check()}, 260 pas par an pour que chaque
date de période tombe sur la grille.
\item \textbf{Réunions}~: \texttt{remaining\_sd\_profile()}~; \texttt{fig\_rc\_rfr.py}
écrit les graphiques.
\end{tutsteps}
\textbf{Pour aller plus loin.} Lancer la simulation avec 250 pas par an et une
grille non alignée (supprimer le contrôle) et observer le biais~; évaluer un caplet sur une période
qui ne contient aucune réunion.
\end{tutorial}

\section{Construction~: caplets rétrospectifs}\label{bld:rc:modelling-overnight-rate-products:rfrcaplet}

\begin{build}
\bfield{Rôle} Caps et floors sur taux overnight composés, les options des
prêts à taux variable du portefeuille depuis la fin des taux interbancaires offerts.
\bfield{Interface} \texttt{hw\_period\_variances}~; \texttt{hw\_backward\_caplet},
\texttt{hw\_forward\_caplet}, \texttt{hw\_backward\_caplet\_mc}~;
\texttt{gfmm\_effective\_variance}, \texttt{black\_caplet}~;
\texttt{meeting\_variance}, \texttt{bachelier\_caplet}.
\bfield{Règles} Capitalisation continue du taux court comme modèle de la
capitalisation quotidienne~; les grilles de simulation doivent contenir chaque date de période (imposé)~; numpy uniquement.
\bfield{Tests d'acceptation} \texttt{code/firm/rfrcaplet/tests/}~: un \omterm{def:rc:modelling-overnight-rate-products:caplet}{caplet rétrospectif}
dépasse son jumeau prospectif de sa variance intra-période, qui vaut
$\sigma^2\delta^3/3$ pour une période commençant maintenant~; la formule fermée concorde avec Monte-Carlo
sur une grille alignée et une grille non alignée est refusée~; la variance à décroissance linéaire et la
variance des réunions admettent leurs formules fermées.
\bfield{Pour aller plus loin} Capitalisation quotidienne avec calendriers de jours ouvrés~; lookback et
délai de paiement~; un \omterm{def:rc:short-rate-models:short}{modèle de taux court} à dates de réunion calibré sur les contrats à terme et leurs options.
\end{build}

\begin{omsources}
\item A.\ Lyashenko and F.\ Mercurio, «~Looking forward to backward-looking rates:
a modeling framework for term rates replacing LIBOR», 2019.
\item O.\ Cheyette, «~Term structure dynamics and mortgage valuation~», \emph{Journal
of Fixed Income} 1(4), 1992.
\item ARRC, «~ARRC formally recommends term SOFR~», 29 juillet 2021.
\item Federal Reserve, \emph{Meeting calendars and information}.
\end{omsources}

\section{Exercices}

\begin{exercise}[$\star$]\label{exo:rc:modelling-overnight-rate-products:1}
Pourquoi un cap sur le SOFR composé inclut-il un caplet sur le trimestre en cours, alors
qu'un cap interbancaire ne l'a jamais fait~?
\end{exercise}

\begin{exercise}[$\star$]\label{exo:rc:modelling-overnight-rate-products:2}
Avec une décroissance linéaire, de combien la variance totale d'un \omterm{def:rc:modelling-overnight-rate-products:caplet}{caplet rétrospectif}
sur $[2,2.25]$ est-elle supérieure à celle de son jumeau prospectif~?
\end{exercise}

\begin{exercise}[$\star$]\label{exo:rc:modelling-overnight-rate-products:3}
Dans l'\cref{ex:rc:modelling-overnight-rate-products:meetings}, calculer les poids
de 69~\% et 23~\% à partir des dates de réunion.
\end{exercise}

\begin{exercise}[$\star\star$]\label{exo:rc:modelling-overnight-rate-products:4}
En utilisant $V_{\text{pendant}}\approx\sigma^2\delta^3/3$ et $V_{\text{avant}}\approx\sigma^2\delta^2S$,
estimer la part intra-période de la variance du deuxième trimestre et la comparer aux
25~\% du texte.
\end{exercise}

\begin{exercise}[$\star\star$]\label{exo:rc:modelling-overnight-rate-products:5}
Lire la \cref{fig:rc:modelling-overnight-rate-products:profile}. À quels jours un caplet court sur le trimestre
perd-il l'essentiel de sa valeur temps selon la vue par réunion, et
qu'est-ce que cela implique pour son theta~?
\end{exercise}

\begin{exercise}[$\star\star$]\label{exo:rc:modelling-overnight-rate-products:6}
Pourquoi les bonnes pratiques restreignent-elles les dérivés sur SOFR à terme, et quelle base
porte un desk qui vend des caps sur taux à terme~?
\end{exercise}

\begin{exercise}[$\star\star\star$]\label{exo:rc:modelling-overnight-rate-products:7}
\emph{Programmation.} Avec \texttt{cap\_table()}, répartir l'écart de 38\,222~USD entre
les deux caps en la part du premier caplet et celle de la variance intra-période des autres.
\end{exercise}

\begin{exercise}[$\star\star\star$]\label{exo:rc:modelling-overnight-rate-products:8}
\emph{Trouver la faille.} «~Nous avons fait passer notre pricer de cap du LIBOR au SOFR en
changeant simplement le nom de l'indice~: il saute toujours le premier caplet et utilise
la variance jusqu'au début de chaque période.~»
\end{exercise}

\section{Problème~: le cap SOFR sur un prêt à effet de levier}

\begin{problem}[{Problème du week-end --- un emprunteur choisit son cap}]\label{pb:rc:modelling-overnight-rate-products:1}
Une entreprise détenue par un fonds de capital-investissement a un prêt de 100~millions d'USD sur deux ans, rémunéré au
SOFR composé à trois mois plus une marge~; ses prêteurs exigent un cap à 3.75~\%. Sa banque propose un
cap sur le SOFR composé et un cap sur le SOFR à terme à trois mois. Utiliser la courbe du chapitre~1 et
le \omterm{def:rc:short-rate-models:hw}{modèle de Hull--White} avec $\kappa=3\%$ et $\sigma=90$ points de base.

\medskip\noindent\textbf{Partie I --- Les deux caps.}
\begin{enumerate}
\item Donner le prix de chaque cap.
\item Pourquoi le cap sur le SOFR composé est-il plus cher~?
\item Donner la valeur du caplet de la période en cours.
\item Quel cap couvre exactement le prêt, et que laisse non couvert l'autre~?
\item Donner la part de variance intra-période pour le deuxième et le huitième trimestre.
\end{enumerate}

\medskip\noindent\textbf{Partie II --- Contrôles.}
\begin{enumerate}[resume]
\item Donner le caplet du cinquième trimestre par la formule fermée et par simulation.
\item Qu'est-ce qui a mal tourné lorsque la simulation a utilisé 250 pas par an~?
\item Que prédit la décroissance linéaire du \omterm{def:rc:modelling-overnight-rate-products:gfmm}{modèle de marché forward généralisé} pour la
variance supplémentaire du cinquième trimestre, et comment se compare-t-elle à celle de Hull--White~?
\item Pourquoi les deux modèles donnent-ils presque la même réponse~?
\item Quels paramètres comptent le plus pour le prix~?
\end{enumerate}

\medskip\noindent\textbf{Partie III --- Réunions.}
\begin{enumerate}[resume]
\item Quelles réunions tombent au cinquième trimestre~?
\item Quand son taux est-il connu dans la vision par réunions~?
\item Quel est le vega du caplet le lendemain de la réunion de décembre~?
\item Comment la banque couvrirait-elle le caplet réunion par réunion~?
\item La vision par réunions modifie-t-elle le prix ou le risque~?
\end{enumerate}

\medskip\noindent\textbf{Partie IV --- Jugement.}
\begin{enumerate}[resume]
\item Quel cap conviendrait-il à l'entreprise d'acheter~?
\item Quel risque de base la banque prend-elle si elle vend le cap sur taux à terme~?
\item Pourquoi les prêteurs acceptent-ils l'un ou l'autre~?
\item Énoncer le \emph{résultat nommé}~: les deux prix de cap et l'écart, avec sa décomposition.
\item En une phrase~: qu'ajoute la mention «~a posteriori~» à une option~?
\end{enumerate}
\end{problem}

\section{Questions d'entretien}

\begin{interviewq}[$\star$ \iqroles{trader, banque}]\label{iq:rc:modelling-overnight-rate-products:1}
Quelle est la différence entre un caplet sur l'Euribor et un caplet sur le taux
€STR composé~?
\end{interviewq}

\begin{interviewq}[$\star\star$ \iqroles{chercheur}]\label{iq:rc:modelling-overnight-rate-products:2}
Comment le modèle de marché doit-il évoluer pour les \omterm{def:rc:modelling-overnight-rate-products:rates}{taux rétrospectifs}~?
\end{interviewq}

\begin{interviewq}[$\star\star$ \iqroles{chercheur, développeur}]\label{iq:rc:modelling-overnight-rate-products:3}
Pricer un \omterm{def:rc:modelling-overnight-rate-products:caplet}{caplet rétrospectif} dans Hull--White. D'où vient la variance
supplémentaire~?
\end{interviewq}

\begin{interviewq}[$\star\star$ \iqroles{trader}]\label{iq:rc:modelling-overnight-rate-products:4}
Comment les réunions des banques centrales modifient-elles le risque des options sur taux au jour le jour à courte échéance~?
\end{interviewq}

\begin{interviewq}[$\star\star\star$ \iqroles{développeur, risque}]\label{iq:rc:modelling-overnight-rate-products:5}
Votre prix Monte Carlo d'un caplet sur taux composé est inférieur de 6~\% à la
formule fermée, avec une erreur type serrée. Que vérifiez-vous en premier~?
\end{interviewq}

\begin{interviewq}[$\star\star\star$ \iqroles{chercheur, banque}]\label{iq:rc:modelling-overnight-rate-products:6}
Pourquoi le SOFR à terme pose-t-il problème comme sous-jacent de dérivés, et que font les banques à la place~?
\end{interviewq}
